% /root/monodromy, 19 November 2026; independently audited by /root.
\subsection{Two exact full monodromy groups on a product surface}
\label{subsec:product-full-groups}
\status{Proved; independent parent proof, primary-source audit and exact-program rerun passed.}

The following computation goes beyond determining invariant covectors.
It determines the entire images in two degrees for one fixed embedded
ordinary abelian surface.  Its use of a projective parameter space is
essential: it permits a short, checkable certificate for all divisorial
inertia, including divisors whose general point parametrizes a singular
curve.

\begin{lemma}[A line certificate for generic stable defect one]
\label{lem:hasse-line-certificate}
Let $M(s)$ be a $b\times b$ matrix of homogeneous linear forms over
$\mathbf F_2$ in homogeneous coordinates on $P=\mathbf P^m$.
Suppose $h(s)=\det M(s)$ is nonzero.  Let $a,b_0$ be coefficient
vectors with $\det M(b_0)=1$.  If, at every root $\alpha$ of
$\det M(a+t b_0)$ over $\overline{\mathbf F}_2$, one has
\[
 \rank\bigl(M(a+\alpha b_0)^{[2]}M(a+\alpha b_0)\bigr)=b-1,
\]
then at the geometric generic point of every irreducible component
of $h=0$ the semilinear map with matrix $M(s)^t$ has stable rank
$b-1$.
\end{lemma}

\begin{proof}
The determinant is homogeneous of degree $b$, and its restriction to
the indicated affine line has degree exactly $b$.  For every geometric
irreducible homogeneous factor $q$ of $h$, one has $q(b_0)\ne0$.
Consequently $q(a+t b_0)$ has its full positive degree and has a root
over the algebraic closure.  If the rank of $M^{[2]}M$ were below
$b-1$ generically on $q=0$, then $q$ would divide every $(b-1)$-minor
of that matrix.  All these minors would vanish at every root of
$q(a+t b_0)$, contrary to the hypothesis.  Thus the first and second
Frobenius images both have dimension $b-1$ at that generic point.
Over a perfect geometric field their equality makes this image
Frobenius-stable and bijective.  This is the asserted stable rank.
The argument does not require the determinant divisor to be reduced.
\end{proof}

\begin{lemma}[Rank-one loss and transvection generation]
\label{lem:char2-defect-one-inertia}
Let $\mathcal E$ be a vector bundle of rank $r$ on a projective space
$P$ over an algebraically closed field of characteristic two, with
a Frobenius-linear endomorphism.  Suppose this map is generically
bijective, and at the generic point of every component of its
noninvertibility divisor its stable rank is $r-1$.
Then the geometric image of the Lang local system on the
invertibility open is generated by transvections.
\end{lemma}

\begin{proof}
At such a divisor, pass to the strict henselization of the discrete
valuation ring.  The Lang scheme $x-F(x)=0$ is separated, affine,
quasi-finite and etale: its Jacobian is the identity, and its
geometric fibres are finite.  Its finite etale part has rank
$2^{r-1}$ and is a subgroup, by
Lemma~\ref{lem:boundary-Lang-finite-part}, or directly by
\cite[Lemma~41.20.5]{Stacks0BTN} and
\cite[Lemma~58.8.2]{Stacks09ZS}.
Since the local ring is strictly henselian, inertia fixes its
$(r-1)$-dimensional generic fibre pointwise.  The quotient has
dimension one over $\mathbf F_2$, so inertia also acts trivially on
that quotient.  Every nonidentity inertia element therefore has
the form $1+u\otimes\lambda$, where $u$ lies in the fixed hyperplane
and $\lambda$ annihilates it.  Its difference from the identity has
rank one and square zero: it is a transvection.

Let $G$ be the geometric image and divide it by the normal subgroup
generated by these divisorial inertia groups.  Normalize $P$ in the
resulting finite Galois extension.  This finite normal cover is
unramified in codimension one and hence is etale everywhere by
purity, \cite[Lemma~58.21.4]{Stacks0BMB}.  Projective space over an
algebraically closed field has no nontrivial connected finite etale
cover, so the quotient is trivial.

For clarity, the last fact can also be proved without a fundamental
group calculation.  For $\mathbf P^1$ it is Riemann--Hurwitz.
For higher-dimensional projective space, pull a cover back to the
blowup of a point.  This blowup is a $\mathbf P^1$-bundle over
$\mathbf P^{m-1}$.  On each fibre the pulled-back cover is trivial;
cohomology and base change applied to its structure sheaf identifies
it with the pullback of a finite etale cover of $\mathbf P^{m-1}$.
The exceptional section is mapped to the original point, so the
restriction of the cover to that section is trivial.  Thus the
cover of $\mathbf P^{m-1}$, and then the pulled-back cover, is
trivial.  The blowup is birational, so the original connected
function-field extension also has degree one.
\end{proof}

\begin{lemma}[A small Frobenius certificate for a transvection group]
\label{lem:product-frobenius-group-test}
Let $N$ be either $8$ or $32$, and let
$H\triangleleft G\subseteq\GL_N(\mathbf F_2)$.
Suppose $H$ is nontrivial and generated by transvections, and
$G$ contains an element whose characteristic polynomial is
irreducible of degree $N$ and is not reciprocal.  Then
$H=G=\GL_N(\mathbf F_2)$.
\end{lemma}

\begin{proof}
The irreducible polynomial makes $G$ irreducible.  The socle of its
restriction to $H$ is nonzero and $G$-stable, so it is the whole
space.  The isotypic $H$-summands are permuted transitively by $G$.
If a simple constituent occurred with multiplicity $e$, the rank
of $h-1$ on its isotypic summand would be a multiple of $e$.
A nontrivial generating transvection has total rank one, so
$e=1$ on one summand and, by transitivity, on all of them.
Thus the space is a direct sum of simple $H$-modules of a common
dimension $c$, and $G$ embeds in
$\GL_c(\mathbf F_2)\wr S_{N/c}$ if there is more than one summand.
This is precisely the elementary argument of
Lemma~\ref{lem:transvection-clifford-blocks}; it includes the
possibility of imprimitive groups rather than assuming it away.

For $N=8$, an irreducible element has order divisible by $17$:
otherwise its order divides $15=2^4-1$, contradicting its degree.
For $N=32$, its order is divisible by $65537$: otherwise it divides
$65535=2^{16}-1$.  Here $65537$ is prime, as direct trial division
up to $257$ verifies, and
\[
 \operatorname{ord}_{17}(2)=8,\qquad
 \operatorname{ord}_{65537}(2)=32.
\]
The indicated prime does not divide $|\GL_c(\mathbf F_2)|$ for
$c<N$, nor $|S_{N/c}|$.  It therefore cannot divide the order of
the displayed proper wreath product.  Thus $H$ is irreducible.

We now use exactly one external finite-group classification:
\cite[Theorem~3.6, printed p.~22]{KantorSeress2001}.
Over $\mathbf F_2$, an irreducible group of dimension at least three
generated by transvection groups is the full special linear group,
or a symplectic group, or one of the listed orthogonal and symmetric
groups in their even-dimensional natural representations.
Each non-linear alternative preserves a nondegenerate alternating
bilinear form.  This form is unique: the endomorphism ring of the
irreducible $H$-module is a finite field, whose degree over
$\mathbf F_2$ divides the rank one of a transvection, and hence is
one.  Schur's lemma then gives a one-dimensional space of invariant
forms.  Every normalizing element of $G$ preserves that form, since
$\mathbf F_2^\times=1$.  Its characteristic polynomial would be
reciprocal, contrary to the hypothesis.  Only
$H=\mathrm{SL}_N(\mathbf F_2)=\GL_N(\mathbf F_2)$ remains.
\end{proof}

\begin{theorem}[Exact groups in degrees one and two on one fixed surface]
\label{thm:product-char2-two-full-groups}
Let
\[
 E:\quad Y^2Z+XYZ+X^3+Z^3=0\quad\text{over }\mathbf F_2,
 \qquad X=E\times E,
 \qquad L=\mathcal O_E(3O)\boxtimes\mathcal O_E(3O),
\]
with the fixed Segre embedding of the two plane cubics.
For the relative Hasse--Witt local system on the ordinary smooth
degree-$d$ hypersurface-section open, both geometric and arithmetic
images are, up to a common choice of basis,
\[
 \begin{array}{c|c|c}
 d&\rank\mathcal V_d&G_{\rm geom}=G_{\rm arith}\\ \hline
 1&8&\GL_8(\mathbf F_2)\\
 2&35&\left\{
       \begin{pmatrix}A&B\\0&I_3\end{pmatrix}:
       A\in\GL_{32}(\mathbf F_2),\ B\in M_{32\times3}(\mathbf F_2)
       \right\}.
 \end{array}
\]
In particular the second image is the full pointwise stabilizer of
the three independent covectors already constructed, of order
$2^{96}|\GL_{32}(\mathbf F_2)|$.  The same statements hold after
base change to any perfect field of characteristic two.
\end{theorem}

\begin{proof}
The generic ranks and the equality between the ambient hypersurface
family and the complete linear system have already been established
in Theorem~\ref{thm:product-exact-fixed-covectors}.  We recall the
matrix model, to make the finite computations reproducible.
Write $x=X/Z$, $y=Y/Z$ and $\omega=dx/x$, normalized by
$C(\omega)=\omega$.  Put $C(f\omega)=T(f)\omega$.
The equation $y^2+xy=x^3+1$ gives
\[
 T(x^{2a})=x^a,\quad T(x^{2a+1})=0,\quad
 T(x^{2a}y)=x^{a+1},\quad
 T(x^{2a+1}y)=x^a(y+1).
\]
For $d=1$ use the ordered basis $(1,x,y)$ of
$H^0(E,\mathcal O(3O))$; for $d=2$ use
$(1,x,x^2,x^3,y,xy)$ of $H^0(E,\mathcal O(6O))$.
Use the lexicographically ordered tensor basis on $X$.
An integer $s=\sum s_j2^j$ encodes the section with coefficient
$s_j\in\mathbf F_2$ on the $j$th tensor basis vector, beginning
with index zero.  Write $W=H^0(X,L^d)$ and $b=\dim W=9d^2$.

For a universal coefficient vector $s$, let $M(s)$ be the linear
matrix whose $j$th column is the coefficient vector of
$T\otimes T(s e_j)$, with the scalar Frobenius root removed.
Thus the actual Cartier map is
$v\mapsto M(s)^{[1/2]}v^{[1/2]}$, and the dual coherent
Frobenius matrix is $M(s)^t$.
The full coherent bundle is $W^*\otimes\mathcal O(-1)$, and its
kernel under evaluation at $s$ is the relative bundle
$\mathcal E_d=\Omega^1_{\mathbf P(W)}$.
Its quotient is $H^2(X,\mathcal O_X)\otimes\mathcal O$, with
bijective Frobenius.  The identity
\[
                       M(s)s=s^{[2]}
\]
expresses this quotient.  Therefore $\det M(s)$ has exactly the
same invertibility open as the relative Frobenius.

This open is defined on the whole projective parameter space,
including singular sections.  The Lang local system extends over
all of it.  Removing the smooth-discriminant locus does not change
its image: a dense open in a normal integral scheme surjects onto
its etale fundamental group.  In particular every ordinary
$\mathbf F_2$ coefficient vector below supplies an actual
arithmetic Frobenius class for the image in question; it need not
itself define a smooth curve.  On the dual representation its
characteristic polynomial is that of the induced Cartier matrix;
on the original representation it is reciprocal.  Our irreducibility,
nonreciprocity, order and fixed-space tests are unchanged by this
duality convention.

\paragraph{Divisorial inertia certificates.}
The retained exact program
\path{code/product_char2_group_certificates.py}
constructs all matrices from the four displayed Cartier formulas.
Its full output, including every binary column and every finite-field
rank, is
\path{outputs/product_char2_group_certificates.json}.
The following two affine lines have ordinary leading coefficient:
\[
 \begin{array}{c|r|r|l}
 d&a&b_0&\det M(a+t b_0)\\ \hline
 1&265&303&t^2(t^2+t+1)(t^5+t^4+t^2+t+1)\\
 2&56866400210&54262944944&(t+1)^4q_7(t)q_9(t)q_{16}(t),
 \end{array}
\]
where
\[
\begin{split}
 q_7&=t^7+t^6+t^5+t^3+t^2+t+1,\\
 q_9&=t^9+t^8+t^6+t^3+1,\\
 q_{16}&=t^{16}+t^{15}+t^{13}+t^{10}+t^6+t^3+t^2+t+1.
\end{split}
\]
The listed factors are irreducible over $\mathbf F_2$.
At a root of each listed factor, both $M$ and $M^{[2]}M$
have rank $8$ for $d=1$, respectively $35$ for $d=2$.
The program verifies irreducibility by the finite-field Frobenius
criterion, multiplies the factors back to the determinant, and
performs Gaussian elimination in the respective finite fields.
The determinants have full degrees $9$ and $36$.
Lemma~\ref{lem:hasse-line-certificate} therefore gives stable rank
$b-1$ for the full coherent module at every geometric generic
Hasse divisor.  Removing its bijective rank-one quotient gives
stable relative rank $b-2$, a loss of exactly one.
Lemma~\ref{lem:char2-defect-one-inertia} proves that both geometric
images are generated by transvections.  This argument includes
all Hasse components without asserting their irreducibility or
reducedness.

\paragraph{Degree one.}
For the ordinary section encoded by $258$, the matrix on $W$
has characteristic polynomial
\[
 (T+1)(T^8+T^4+T^3+T+1).
\]
The degree-eight factor is irreducible and not reciprocal.
For the ordinary section encoded by $27$, the characteristic
polynomial is $(T+1)^9$.
Thus the geometric image is not trivial: if it were trivial,
the local system would be constant over the algebraic closure,
and all rational points would have the same arithmetic Frobenius
characteristic polynomial.  The arithmetic image is irreducible
because of the first section.  Apply
Lemma~\ref{lem:product-frobenius-group-test} with $N=8$.
It gives both images equal to $\GL_8(\mathbf F_2)$.

\paragraph{The action on the common kernel in degree two.}
The four character projections $s_\chi$ of $s$ give four fixed
vectors in $W$, with sum $s$; in $W/\langle s\rangle$ they give
three independent invariant covectors.  All are defined over
$\mathbf F_2$, by the explicit Cartier calculation above and
Theorem~\ref{thm:product-exact-fixed-covectors}.
Let $K\subset V$ be their common kernel.  Then $\dim K=32$ and
the quotient $Q=V/K$ is the constant three-dimensional module
for both geometric and arithmetic monodromy.

At the ordinary section encoded by $17707720254$, the characteristic
polynomial on $W$ is $(T+1)^4f_{32}(T)$, where
\[
\begin{split}
 f_{32}(T)={}&T^{32}+T^{28}+T^{27}+T^{26}+T^{25}+T^{24}
             +T^{23}+T^{22}+T^{20}+T^{17}\\
            &+T^{14}+T^{11}+T^8+T^6+T^5+T^3+T^2+T+1.
\end{split}
\]
The program verifies that this polynomial is irreducible and not
reciprocal.  Consequently the arithmetic action on $K$ is
irreducible.  At the second ordinary section, encoded by
$65866883812$, the polynomial on $K$ has a root at $1$.
Thus the geometric image on $K$ is nontrivial, by the same
constant-local-system argument as in degree one.

The restriction to $K$ of a transvection of $V$ is either the
identity or a transvection: every image line lies in $K$, since
the three quotient covectors are fixed.  Hence the geometric
image on $K$ is generated by transvections.  Applying
Lemma~\ref{lem:product-frobenius-group-test} with $N=32$ shows
that this image is the entire $\GL(K)$.

\paragraph{The translation subgroup in degree two.}
It remains to prove that every map $Q\to K$ occurs; surjectivity
onto $\GL(K)$ alone does not prove this.
Choose a vector-space splitting $V=K\oplus Q$ and let $G_0$ be
the geometric image.  Its kernel $R$ in
\[
 1\longrightarrow R\longrightarrow G_0
          \longrightarrow\GL(K)\longrightarrow1
\]
is a $\GL(K)$-submodule of $\Hom(Q,K)=K\otimes Q^*$.
The natural module is absolutely simple with endomorphism ring
$\mathbf F_2$.  Thus $R=K\otimes U$ for a subspace $U\subseteq Q^*$.
Modulo $R$, the subgroup is the graph of a cocycle in
$H^1(\GL(K),K\otimes(Q^*/U))$.
This cohomology is zero by the elementary transvection calculation
in Lemma~\ref{lem:char2-group-inputs}, valid in every dimension
at least four.  Conjugating the splitting therefore makes
$G_0=R\rtimes\GL(K)$, and its invariant-vector space projects
isomorphically onto $U^\perp\subseteq Q$.

That invariant-vector space is stable under arithmetic monodromy.
Its projection to the constant quotient $Q$ is injective, so
arithmetic monodromy acts trivially on it as well.
We now use the additional exact certificate at $s=65866883812$:
if $C=M(s)$ and the four character vectors are $s_\chi$, then
\begin{equation}\label{eq:product-no-complement-certificate}
 \dim\bigl(\im(C-I)+\langle s\rangle\bigr)
 =\dim\bigl(\im(C-I)+\langle s,s_\chi:\chi\rangle\bigr)=33.
\end{equation}
The full columns and the four vectors are retained in the output.
On the relative coherent fibre $V\otimes\bar k=\operatorname{ann}(s)$,
a vector fixed by this rational-point Frobenius lies in
$\ker(C^t-I)\cap\operatorname{ann}(s)$.
Equation~\eqref{eq:product-no-complement-certificate} forces it
to annihilate every $s_\chi$ and hence to lie in $K$.
Every geometric invariant vector is arithmetic invariant, so this
applies to it.  But $K^{\GL(K)}=0$.  Therefore $V^{G_0}=0$,
$U^\perp=0$, and $U=Q^*$.
The kernel is the whole $\Hom(Q,K)$, proving the asserted full
pointwise stabilizer.  The arithmetic image is contained in that
same stabilizer and contains the geometric image, so the two agree.

Finally the geometric covers remain connected under algebraically
closed extension.  Both displayed groups are already defined over
$\mathbf F_2$, and the three degree-two covectors are rational
over that field.  After any perfect base extension the arithmetic
image is therefore still between the same equal geometric image
and the same upper bound.  This proves the final assertion.
\end{proof}

\begin{remark}[Verification scope]
The group conclusions use the lemmas proved here, the elementary
cocycle and socle arguments already proved, and the single opened
finite-group classification specified above.  The retained program
provides exact polynomial, rank and Frobenius certificates.  It does
not infer monodromy from a random sample.  The exploratory search
chooses certificates, after which every statement needed from each
certificate is verified exactly.  The high-dimensional characteristic
polynomials are also recomputed independently by SymPy's integer
characteristic-polynomial routine and reduced modulo two.
No assertion for degree $d\ge3$ on this surface follows just from
these two degree calculations.
\end{remark}
